\documentclass[letterpaper]{article} % DO NOT CHANGE THIS
\usepackage{aaai24}  % DO NOT CHANGE THIS
\usepackage{times}  % DO NOT CHANGE THIS
\usepackage{helvet}  % DO NOT CHANGE THIS
\usepackage{courier}  % DO NOT CHANGE THIS
\usepackage[hyphens]{url}  % DO NOT CHANGE THIS
\usepackage{graphicx} % DO NOT CHANGE THIS
\urlstyle{rm} % DO NOT CHANGE THIS
\def\UrlFont{\rm}  % DO NOT CHANGE THIS
\usepackage{natbib}  % DO NOT CHANGE THIS AND DO NOT ADD ANY OPTIONS TO IT
\usepackage{caption} % DO NOT CHANGE THIS AND DO NOT ADD ANY OPTIONS TO IT
\frenchspacing  % DO NOT CHANGE THIS
\setlength{\pdfpagewidth}{8.5in}  % DO NOT CHANGE THIS
\setlength{\pdfpageheight}{11in}  % DO NOT CHANGE THIS




\usepackage{fancyhdr}

\fancyhf{} % sets both header and footer to nothing
\renewcommand{\headrulewidth}{0pt}
\fancyhf{} % sets both header and footer to nothing
\setlength{\headheight}{12pt}
\chead{\fontfamily{phv}\selectfont \small This paper has been accepted for publication at the AAAI 2024 Workshop on AI to Accelerate Science and Engineering}

\usepackage{algorithm}
\usepackage{algorithmic}

\usepackage{newfloat}
\usepackage{listings}
\DeclareCaptionStyle{ruled}{labelfont=normalfont,labelsep=colon,strut=off} % DO NOT CHANGE THIS
\lstset{%
	basicstyle={\footnotesize\ttfamily},% footnotesize acceptable for monospace
	numbers=left,numberstyle=\footnotesize,xleftmargin=2em,% show line numbers, remove this entire line if you don't want the numbers.
	aboveskip=0pt,belowskip=0pt,%
	showstringspaces=false,tabsize=2,breaklines=true}
\floatstyle{ruled}
\newfloat{listing}{tb}{lst}{}
\floatname{listing}{Listing}
\usepackage{subfiles}
\usepackage{xcolor}
\setcounter{secnumdepth}{0} %May be changed to 1 or 2 if section numbers are desired.


\title{Position Paper on Materials Design -- A Modern Approach}
\author{
    Willi Gro{\ss}mann\textsuperscript{\rm 1}, %\equalcontrib,
    Sebastian Eilermann\textsuperscript{\rm 1}, %\equalcontrib,
    Tim Rensmeyer\textsuperscript{\rm 1}, \\
    Artur Liebert\textsuperscript{\rm 1},  Michael Hohmann\textsuperscript{\rm 1}, Christian Wittke\textsuperscript{\rm 1},
    Oliver Niggemann\textsuperscript{\rm 1}
}
\affiliations{
\textsuperscript{\rm 1}Institute of Automation Technology\\Helmut-Schmidt-Universitt / Universitt der Bundeswehr Hamburg, \\
Holstenhofweg 85, 22043 Hamburg, Germany\\
}

\iffalse
\title{My Publication Title --- Single Author}
\author {
    Author Name
}
\affiliations{
    Affiliation\\
    Affiliation Line 2\\
    name@example.com
}
\fi

\iffalse
\title{My Publication Title --- Multiple Authors}
\author {
    First Author Name\textsuperscript{\rm 1,\rm 2},
    Second Author Name\textsuperscript{\rm 2},
    Third Author Name\textsuperscript{\rm 1}
}
\affiliations {
    \textsuperscript{\rm 1}Affiliation 1\\
    \textsuperscript{\rm 2}Affiliation 2\\
    firstAuthor@affiliation1.com, secondAuthor@affilation2.com, thirdAuthor@affiliation1.com
}
\fi
\usepackage{bibentry}

\begin{document}

\maketitle
\thispagestyle{fancy}

\begin{abstract}
Traditional design cycles for new materials and assemblies have two fundamental drawbacks. The underlying physical relationships are often too complex to be precisely calculated and described. Aside from that, many unknown uncertainties, such as exact manufacturing parameters or materials composition, dominate the real assembly behavior. Machine learning (ML) methods overcome these fundamental limitations through data-driven learning. In addition, modern approaches can specifically increase system knowledge. Representation Learning allows the physical, and if necessary, even symbolic interpretation of the learned solution. In this way, the most complex physical relationships can be considered and quickly described. Furthermore, generative ML approaches can synthesize possible morphologies of the materials based on defined conditions to visualize the effects of uncertainties. This modern approach accelerates the design process for new materials and enables the prediction and interpretation of realistic materials behavior.

\end{abstract}

\subfile{sections/1_Intro}

\subfile{sections/2_SoA}

\subfile{sections/3_Solution}

\subfile{sections/4_Conclusion}

\section{Acknowledgment}
\label{Ackn}
This work was supported by the project "KIBIDZ --Intelligente Brandgefahrenanalyse f\"ur Geb\"aude und Schutz der Rettungskrafte durch K\"unstliche Intelligenz und Digitale Brandgeb\"audezwillinge" by the Zentrum f\"ur Digitalisierungs- und Technologieforschung der Bundeswehr (dtec.bw).


Odio architecto voluptas cumque, adipisci id aspernatur consectetur porro cumque nam quas officia repudiandae facilis, sed ullam minima soluta quod quis porro, distinctio dolor illum provident repellat natus inventore?Consequuntur corporis explicabo alias debitis perspiciatis possimus nihil nulla vitae, aliquid dolorem dolore at ipsa assumenda, nam explicabo laboriosam impedit culpa, repellendus laborum sapiente sit quod expedita laudantium illo quaerat.Illum eligendi facilis odio quam, inventore ducimus vel doloribus voluptate, voluptates quod quaerat ipsum officiis sed, perspiciatis facere facilis doloremque, repellat nobis eligendi quos nemo quidem ullam in possimus laboriosam.Quam veniam maxime dignissimos architecto minima autem molestias, delectus omnis nostrum sequi sunt itaque ea reprehenderit, neque at error consequatur eos consectetur ullam et animi, beatae repellendus dolorem asperiores vero eos aliquid ratione quas recusandae exercitationem minima?Consectetur natus hic aliquam dolore error suscipit eos recusandae, incidunt corporis at quas voluptatibus quibusdam accusantium officia, nam obcaecati doloribus voluptatibus dolores natus culpa eos.Reiciendis quaerat maiores praesentium quo minus voluptate eaque qui minima, voluptatem laborum earum.Ullam error quae facere labore, iure id suscipit fuga dolorem repudiandae obcaecati molestias excepturi.Repellat velit numquam pariatur odit nam ipsum tempore, minus optio harum beatae totam eos
\bibliography{lit}

\end{document}